\documentclass[12pt]{article}

\usepackage[margin=1in]{geometry}
\usepackage{booktabs}
\usepackage{array}
\usepackage{tabularx}
\usepackage{hyperref}
\usepackage{amsmath}
\usepackage{amssymb}
\usepackage{natbib}
\usepackage{setspace}
\usepackage{parskip}

\onehalfspacing

\title{Treatment of Error Covariance, Clustering, and Serial Correlation\\
       in the Staggered Difference-in-Differences Literature:\\
       A Survey of Six Key Papers}
\author{Arora \& Bijani (2026)\\
        \small{Internal Research Note}}
\date{\today}

\begin{document}

\maketitle
\tableofcontents
\newpage

% ============================================================
\section{Introduction}
% ============================================================

A central but often understated design choice in difference-in-differences (DiD)
estimation is how to handle the error covariance matrix $\Sigma$.
The validity of inference and the efficiency of estimation both depend critically
on (i) what structure is imposed on $\Sigma$ and (ii) how that structure is
exploited or approximated in the standard-error formula.

This note surveys six influential papers in the staggered DiD literature through
the lens of error covariance treatment:

\begin{enumerate}
  \item \citet{callaway2021difference} --- Difference-in-differences with multiple
        time periods.
  \item \citet{dechaisemartin2023twoway} --- Two-way fixed effects and
        difference-in-differences with heterogeneous treatment effects: a survey.
  \item \citet{gardner2024twostage} --- Two-stage differences in differences.
  \item \citet{sun2021estimating} --- Estimating dynamic treatment effects in event
        studies with heterogeneous treatment effects.
  \item \citet{wooldridge2025twoway} --- Two-way fixed effects, the two-way Mundlak
        regression, and difference-in-differences estimators.
  \item \citet{borusyak2024revisiting} --- Revisiting event-study designs: robust
        and efficient estimation.
\end{enumerate}

For each paper we address four questions:
\begin{itemize}
  \item What structure is placed on the error covariance (diagonal, block-diagonal,
        fully general)?
  \item What variance estimator is recommended (cluster-robust, bootstrap,
        analytical)?
  \item Is an efficient (GLS-type) estimator proposed that exploits $\Sigma$?
  \item What caveats or special cases are highlighted (few clusters,
        cross-sectional dependence, heteroskedasticity)?
\end{itemize}

% ============================================================
\section{Paper-by-Paper Summaries}
% ============================================================

% ------------------------------------------------------------
\subsection{Callaway and Sant'Anna (2021)}
% ------------------------------------------------------------

\paragraph{Error structure.}
\citet{callaway2021difference} assume that units $i = 1, \ldots, N$ are drawn
i.i.d.\ from the population.
This imposes \textit{cross-sectional independence}: the $N \times N$ block of the
covariance matrix corresponding to distinct units is zero.
Within a unit, however, the covariance matrix $\Sigma_i$ over the $T$ time periods
is left \textit{completely unrestricted}.
The aggregate covariance matrix is therefore block-diagonal with $N$ blocks of
dimension $T \times T$, where each block is an arbitrary positive semi-definite
matrix.

\paragraph{Variance estimator.}
Standard errors are derived via the \textbf{influence-function (IF) approach}.
For each ATT$(g,t)$ estimand, the authors derive an influence function
$\psi_{it}(g,t)$ at the unit level.
The asymptotic variance is then estimated by averaging the squared IFs across
units, which automatically delivers a \textbf{cluster-robust standard error}
at the unit level and is valid under the arbitrary within-unit serial
correlation permitted by the block-diagonal structure.
For simultaneous (uniform) inference across multiple $(g,t)$ cells, they
recommend the \textbf{multiplier (wild) bootstrap}, which resamples i.i.d.\
multipliers at the unit level and thereby preserves within-unit temporal
dependence.

\paragraph{GLS / efficiency.}
No efficient GLS estimator is proposed.
The doubly-robust estimator achieves the semiparametric efficiency bound for
each ATT$(g,t)$ within the class of estimators satisfying their identification
assumptions, but it does not model or exploit the within-unit covariance
$\Sigma_i$.

\paragraph{Caveats.}
When the number of treated cohorts is small, multiplier bootstrap critical values
can under-cover because bootstrap quantiles of the maximum statistic are poorly
estimated with few groups.
In such cases, the authors recommend clustering at the treatment-cohort level
(more conservative) or using randomisation inference.

% ------------------------------------------------------------
\subsection{De Chaisemartin and d'Haultfoeuille (2023)}
% ------------------------------------------------------------

\paragraph{Error structure.}
As a survey, \citet{dechaisemartin2023twoway} do not impose a single covariance
assumption across all reviewed estimators.
For their own DID$_M$ and related estimators, the maintained assumption is
cross-sectional independence with unrestricted within-group serial correlation
--- the same block-diagonal structure as \citet{callaway2021difference}.
The survey prominently cites \citet{bertrand2004much}, who document that
ignoring within-unit serial correlation in TWFE leads to severely
under-estimated standard errors in long panels, because the pre- and
post-treatment group means are highly correlated within units.

\paragraph{Variance estimator.}
For their DID$_M$ estimator, the authors use \textbf{influence-function-based
standard errors} clustered at the (group, period) level.
For the broader class of estimators reviewed, they recommend clustering at the
highest feasible level (typically the unit or treatment-group level) to guard
against serial correlation.

\paragraph{GLS / efficiency.}
Not addressed.
The paper's focus is identification and bias correction for heterogeneous
treatment effects rather than efficiency.

\paragraph{Caveats.}
The \textbf{few-clusters problem} receives substantial attention.
When the number of treated groups is small (fewer than roughly 20),
cluster-robust standard errors from OLS/WLS have poor finite-sample coverage.
The authors recommend the wild cluster bootstrap or randomisation inference
as alternatives.
They also note that heterogeneous treatment effects can interact with
unbalanced panel structures to amplify the bias from ignoring serial
correlation.

% ------------------------------------------------------------
\subsection{Gardner, Thakral, To and Luther (2024)}
% ------------------------------------------------------------

\paragraph{Error structure.}
\citet{gardner2024twostage} assume cross-sectional independence with arbitrary
within-unit serial correlation (block-diagonal $\Sigma$).
A subtle additional consideration arises from the two-stage structure: the
first-stage regression of outcomes on unit and time fixed effects uses only
the untreated observations.
The resulting residual for the second stage is therefore a function of all
untreated observations, creating a nuisance dependence across units that must
be handled carefully in the variance derivation.

\paragraph{Variance estimator.}
The estimator is cast as a \textbf{GMM estimator}, which delivers a natural
\textbf{cluster-robust sandwich variance} formula that accounts for both stages.
A key result is that the ``naive'' second-stage clustered standard error ---
obtained by simply clustering the OLS residuals from the second stage while
ignoring first-stage uncertainty --- is \textit{asymptotically valid} when the
first-stage model is correctly specified, because the first-stage estimation
error washes out in large samples.
Under misspecification or heterogeneous treatment effects, the full two-stage
sandwich correction is required.

\paragraph{GLS / efficiency.}
No efficient GLS estimator is proposed.
The authors note that efficiency improvements from modelling $\Sigma$ are
possible but leave this to future work.

\paragraph{Caveats.}
The asymptotic validity of the naive second-stage SE relies on correctly
specifying the first-stage potential outcome model.
If the first-stage fixed-effect structure is misspecified (e.g., treatment
anticipation), first-stage uncertainty becomes non-negligible even asymptotically.

% ------------------------------------------------------------
\subsection{Sun and Abraham (2021)}
% ------------------------------------------------------------

\paragraph{Error structure.}
\citet{sun2021estimating} adopt a standard linear panel model with no special
structure on the error covariance beyond what is needed for OLS consistency.
Cross-sectional independence is assumed; within-unit temporal dependence is
permitted but not modelled.

\paragraph{Variance estimator.}
Because their interaction-weighted (IW) estimator is a linear transformation of
OLS coefficients from a regression with cohort $\times$ time dummy interactions,
\textbf{standard cluster-robust standard errors from OLS apply directly}.
No separate asymptotic derivation is needed: the delta-method applied to the
OLS covariance matrix delivers the variance of the weighted average of
cohort-specific ATTs.
Clustering at the unit level handles within-unit serial correlation.

\paragraph{GLS / efficiency.}
Not discussed.
The paper's contribution is the identification of a clean estimand that avoids
contamination from heterogeneous effects in TWFE event studies, rather than
an improvement in efficiency or a new inference procedure.

\paragraph{Caveats.}
Because the IW estimator inherits the inference properties of TWFE OLS,
it shares TWFE's sensitivity to heteroskedasticity and serial correlation
when standard (non-cluster-robust) standard errors are used.
The paper offers no guidance on inference with few clusters.

% ------------------------------------------------------------
\subsection{Wooldridge (2025)}
% ------------------------------------------------------------

\paragraph{Error structure.}
\citet{wooldridge2025twoway} works within the \textit{correlated random effects}
(CRE) / two-way Mundlak (TWM) framework, which adds cohort-period means of
time-varying covariates to absorb level correlations between unit effects and
regressors.
The residual idiosyncratic error $\varepsilon_{it}$ after the Mundlak projection
is allowed to have arbitrary within-unit serial correlation.
For inference via pooled OLS, the covariance is treated as block-diagonal.
For GLS estimation, Wooldridge allows more structured within-unit covariances,
including AR(1) and equicorrelated (random-effects) structures.

\paragraph{Variance estimator.}
For pooled OLS on the extended TWFE (ETWFE) or TWM specification, the
recommendation is \textbf{cluster-robust standard errors at the unit level},
valid for large $N$ regardless of $T$ and robust to arbitrary within-unit serial
dependence.
In large-$T$ contexts where $T / N$ is non-negligible, time-series HAC
corrections may be preferred.

\paragraph{GLS / efficiency.}
This is the most substantive treatment of efficiency among the first five papers.
Wooldridge explicitly derives conditions under which \textbf{feasible GLS (FGLS)}
on the ETWFE specification is more efficient than pooled OLS.
When $\Sigma$ follows an AR(1) or equicorrelated structure and is correctly
specified, FGLS with estimated $\hat{\Sigma}$ achieves asymptotic efficiency gains.
He notes that cohort-time cell sizes are important: when cohort sizes are small,
the ATT estimator from ETWFE has high variance, and efficiency-weighted
aggregation via GLS can be substantially more precise.
Even under misspecification of $\Sigma$, FGLS is consistent; only efficiency
is lost.

\paragraph{Caveats.}
The Mundlak device controls for correlated random effects but does not
fully account for all sources of within-unit serial correlation in
$\varepsilon_{it}$.
Post-Mundlak residuals may still exhibit serial dependence, so clustering
remains necessary even after FGLS.

% ------------------------------------------------------------
\subsection{Borusyak, Jaravel and Spiess (2024)}
% ------------------------------------------------------------

\paragraph{Error structure.}
\citet{borusyak2024revisiting} provide the most explicit and general treatment of
the error covariance among all six papers.
They write the untreated potential outcome as
\[
  Y_{it}(0) = \alpha_i + \lambda_t + \varepsilon_{it},
\]
where the idiosyncratic error $\varepsilon_{it}$ is allowed to have:
\begin{itemize}
  \item \textbf{Arbitrary within-unit serial correlation}: $\mathrm{Var}(\varepsilon_i)
        = \Sigma_i$, a $T \times T$ positive definite matrix specific to unit $i$.
  \item \textbf{Heteroskedasticity across units}: $\Sigma_i$ may differ by unit.
  \item \textbf{Cross-sectional independence}: units are i.i.d.\ draws from the
        population; no spatial, network, or cross-sectional dependence is
        accommodated.
\end{itemize}

\paragraph{Variance estimator.}
Two variance estimators are derived:
\begin{enumerate}
  \item \textbf{Cluster-robust sandwich estimator.}
        Valid under cross-sectional independence with arbitrary $\Sigma_i$.
        This is their ``robust'' standard error and is analogous to
        heteroskedasticity-and-autocorrelation-consistent (HAC) variance
        estimation at the unit (cluster) level.
  \item \textbf{Efficient (GLS-based) sandwich estimator.}
        For the efficient imputation estimator, they propose estimating a common
        $\hat{\Sigma}$ from pre-treatment residuals (assuming $\Sigma_i = \Sigma$
        for all $i$), then constructing a feasible GLS sandwich variance.
        This estimator is consistent even under mild misspecification of $\Sigma$,
        because the sandwich form corrects for model error.
\end{enumerate}

\paragraph{GLS / efficiency.}
This is the paper's central methodological contribution on the inference side.
The authors prove that the \textbf{imputation-based GLS estimator achieves the
semiparametric Cram\'{e}r--Rao lower bound} for the ATT$(g,t)$ under their model.
The argument proceeds as follows:
\begin{enumerate}
  \item The untreated potential outcome model is fitted to control observations
        by GLS using the estimated $\hat{\Sigma}$.
  \item Counterfactual outcomes for treated units are imputed from this model.
  \item The treatment effect is the gap between observed and imputed outcomes,
        averaged with GLS weights.
\end{enumerate}
They show that estimators based on OLS (including
\citealt{callaway2021difference},
\citealt{sun2021estimating}, and
\citealt{gardner2024twostage})
are \textit{not} efficient because they do not exploit the within-unit
covariance structure.
The efficiency gain is substantial when within-unit serial correlation is high
or when unit-level variances are heterogeneous.

\paragraph{Caveats.}
Estimating $\Sigma$ consistently from pre-treatment residuals requires a
sufficiently long pre-treatment window ($T_{\text{pre}} \gg T_{\text{post}}$).
With short panels, the plug-in GLS estimator may be more variable than OLS
in finite samples, eroding the asymptotic efficiency gain.
The efficiency result is established under \textit{homogeneous} $\Sigma$ across
units; with heterogeneous $\Sigma_i$ the GLS estimator remains consistent but
loses the formal Cram\'{e}r--Rao guarantee.
Cross-sectional dependence (spatial, network) is not addressed.

% ============================================================
\section{Comparative Summary}
% ============================================================

Table~\ref{tab:summary} consolidates the four dimensions across all six papers.

\begin{table}[ht]
\centering
\caption{Error Covariance Treatment Across Six DiD Papers}
\label{tab:summary}
\renewcommand{\arraystretch}{1.4}
\begin{tabularx}{\textwidth}{@{}lXXcc@{}}
\toprule
\textbf{Paper} & \textbf{Error covariance assumed} & \textbf{SE recommendation}
  & \textbf{GLS?} & \textbf{Efficiency bound?} \\
\midrule
Callaway \& Sant'Anna (2021)
  & Block-diagonal; unrestricted within-unit $\Sigma_i$
  & Influence-function; multiplier bootstrap for joint inference
  & No & No \\
De Chaisemartin \& d'Haultfoeuille (2023)
  & Block-diagonal; unrestricted within-unit
  & IF-based; wild bootstrap for few clusters
  & No & No \\
Gardner et al.\ (2024)
  & Block-diagonal; unrestricted within-unit
  & GMM cluster-robust sandwich (two-stage corrected)
  & No & No \\
Sun \& Abraham (2021)
  & Block-diagonal; unrestricted within-unit
  & Standard OLS cluster-robust at unit level
  & No & No \\
Wooldridge (2025)
  & Block-diagonal; AR(1) or equicorrelated for FGLS
  & Cluster-robust OLS; FGLS for efficiency
  & Yes & No \\
Borusyak et al.\ (2024)
  & Block-diagonal; unrestricted $\Sigma_i$; heteroskedastic
  & Cluster-robust sandwich; GLS sandwich for efficient version
  & \textbf{Yes} & \textbf{Yes} \\
\bottomrule
\end{tabularx}
\end{table}

% ============================================================
\section{Key Takeaways}
% ============================================================

\begin{enumerate}

  \item \textbf{Universal cross-sectional independence.}
        All six papers assume that units are independent across one another.
        None addresses spatial correlation, network dependence, or other
        forms of cross-sectional dependence.
        This is an important gap: in many empirical applications (regional
        policy, financial markets, network effects) cross-unit dependence
        is plausible and could invalidate standard inference.

  \item \textbf{Universal within-unit flexibility.}
        All six papers allow arbitrary within-unit serial correlation over
        time (block-diagonal $\Sigma$ with unrestricted $T \times T$ blocks).
        Cluster-robust standard errors at the unit level are the consensus
        recommendation for handling this.

  \item \textbf{Only two papers propose GLS.}
        \citet{wooldridge2025twoway} and \citet{borusyak2024revisiting} go
        beyond cluster-robust OLS to propose efficient estimators that
        exploit the within-unit covariance structure.
        Wooldridge's approach requires a parametric model for $\Sigma$
        (AR(1) or equicorrelated); Borusyak et al.\ allow a non-parametric
        $\hat{\Sigma}$ estimated from pre-treatment residuals.

  \item \textbf{Borusyak et al.\ provide the only efficiency bound.}
        Of the six papers, only \citet{borusyak2024revisiting} formally
        derive the semiparametric Cram\'{e}r--Rao lower bound for the
        ATT$(g,t)$ and prove that their imputation GLS estimator achieves it.
        The efficiency gain over OLS-based estimators can be substantial in
        long panels with high serial correlation.

  \item \textbf{The few-clusters problem is underexplored.}
        Only \citet{dechaisemartin2023twoway} and \citet{callaway2021difference}
        discuss the case where the number of treatment clusters is small.
        The other papers implicitly assume a large number of clusters, which
        may not hold in typical DiD applications where treatment is assigned
        at the state, industry, or firm level.

  \item \textbf{Implications for the GMM framework.}
        The GMM framework in \citet{arora2026gmm} explicitly models the
        moment covariance $\Omega_\phi$ using estimated autocovariances
        $\hat{\sigma}_d$ from within-unit residuals.
        This is closest in spirit to the efficient approach of
        \citet{borusyak2024revisiting}: both exploit the panel time-series
        structure of $\Sigma$ to construct efficient weights.
        Unlike Borusyak et al., however, the GMM approach additionally
        handles ``forbidden'' comparisons and imposes the linear restrictions
        of the DiD incidence matrix $Q_H$, making it applicable in designs
        where no clean ``imputation'' control group exists.

\end{enumerate}

% ============================================================
\bibliographystyle{apalike}
\bibliography{references}
% ============================================================

% Inline bibliography for self-contained document
\begin{thebibliography}{99}

\bibitem[Arora and Bijani(2026)]{arora2026gmm}
Arora, P. and Bijani, R. (2026).
\newblock Estimating treatment effects under staggered timing using {GMM}.
\newblock Working paper.

\bibitem[Bertrand et al.(2004)]{bertrand2004much}
Bertrand, M., Duflo, E., and Mullainathan, S. (2004).
\newblock How much should we trust differences-in-differences estimates?
\newblock \textit{Quarterly Journal of Economics}, 119(1):249--275.

\bibitem[Borusyak et al.(2024)]{borusyak2024revisiting}
Borusyak, K., Jaravel, X., and Spiess, J. (2024).
\newblock Revisiting event-study designs: robust and efficient estimation.
\newblock \textit{Review of Economic Studies}, 91(6):3253--3285.

\bibitem[Callaway and Sant'Anna(2021)]{callaway2021difference}
Callaway, B. and Sant'Anna, P.~H. (2021).
\newblock Difference-in-differences with multiple time periods.
\newblock \textit{Journal of Econometrics}, 225(2):200--230.

\bibitem[De~Chaisemartin and d'Haultfoeuille(2023)]{dechaisemartin2023twoway}
De~Chaisemartin, C. and d'Haultfoeuille, X. (2023).
\newblock Two-way fixed effects and differences-in-differences with
  heterogeneous treatment effects: a survey.
\newblock \textit{The Econometrics Journal}, 26(3):C1--C30.

\bibitem[Gardner et al.(2024)]{gardner2024twostage}
Gardner, J., Thakral, N., To, L.~T., and Luther, Y. (2024).
\newblock Two-stage differences in differences.
\newblock Working paper.

\bibitem[Sun and Abraham(2021)]{sun2021estimating}
Sun, L. and Abraham, S. (2021).
\newblock Estimating dynamic treatment effects in event studies with
  heterogeneous treatment effects.
\newblock \textit{Journal of Econometrics}, 225(2):175--199.

\bibitem[Wooldridge(2025)]{wooldridge2025twoway}
Wooldridge, J.~M. (2025).
\newblock Two-way fixed effects, the two-way {Mundlak} regression, and
  difference-in-differences estimators.
\newblock \textit{Empirical Economics}, 69(5):2545--2587.

\end{thebibliography}

\end{document}
